\begin{Solbook}{1}
 \par \noindent  \begin {enumerate} \item Điều kiện xác định: $x\ne 2$, $x\ne -2$. \begin {align*} \dfrac {3x+5}{x-2}-\dfrac {3x-1}{x+2}&=\dfrac {26}{x^2-4}\\ (3x+5)(x+2)-(3x-1)(x-2)&=26\\ 18x+8&=26\\ x&=1. \end {align*} Giá trị $x=1$ thỏa mãn điều kiện xác định. Vậy phương trình có nghiệm $x=1$. \par \item Nhân cả hai vế với $12>0$, ta được \begin {align*} 4(2x-3)-3(5x+1)&\ge 2(1-2x)\\ 8x-12-15x-3&\ge 2-4x\\ -3x&\ge 17\\ x&\le -\dfrac {17}{3}. \end {align*} Vậy bất phương trình có nghiệm $x\le -\dfrac {17}{3}$. \par \item Dùng phương pháp thế, ta có \begin {align*} &\begin {cases} 2(x-1)-3(y+2)=-11,\\ \dfrac {x+2y}{3}+\dfrac {x-y}{2}=3 \end {cases}\\[0.7em] \Leftrightarrow \quad &\begin {cases} 2x-3y=-3,\\ 5x+y=18 \end {cases}\\[0.7em] \Leftrightarrow \quad &\begin {cases} 2x-3(18-5x)=-3,\\ y=18-5x \end {cases}\\[0.7em] \Leftrightarrow \quad &\begin {cases} 17x=51,\\ y=18-5x \end {cases}\\[0.7em] \Leftrightarrow \quad &\begin {cases} x=3,\\ y=3. \end {cases} \end {align*} Vậy hệ phương trình có nghiệm $(x;y)=(3;3)$. \end {enumerate}  \hfill \qedEX 
\end{Solbook}
\begin{Solbook}{2}
 \par \noindent  \begin {enumerate} \setlength {\itemsep }{1em} \par \item Ta có \begin {align*} 14+6\sqrt {5} &=9+2\cdot 3\sqrt {5}+5\\ &=\left (3+\sqrt {5}\right )^2,\\ 17-12\sqrt {2} &=9-2\cdot 3\cdot 2\sqrt {2}+8\\ &=\left (3-2\sqrt {2}\right )^2. \end {align*} Vì $3-2\sqrt {2}>0$ nên \begin {align*} \sqrt {14+6\sqrt {5}}-\sqrt {17-12\sqrt {2}} &=\sqrt {\left (3+\sqrt {5}\right )^2} -\sqrt {\left (3-2\sqrt {2}\right )^2}\\ &=3+\sqrt {5}-\left (3-2\sqrt {2}\right )\\ &=\sqrt {5}+2\sqrt {2}. \end {align*} Vậy biểu thức đã cho rút gọn thành $\sqrt {5}+2\sqrt {2}$. \par \item Ta có $\sqrt {45}=3\sqrt {5}$, do đó \begin {align*} \dfrac {\sqrt {45}+\sqrt {5}}{\sqrt {5}-\sqrt {3}} -\dfrac {4\sqrt {3}}{\sqrt {5}+\sqrt {3}}\\ &=\dfrac {4\sqrt {5}}{\sqrt {5}-\sqrt {3}} -\dfrac {4\sqrt {3}}{\sqrt {5}+\sqrt {3}}\\ &=\dfrac {4\sqrt {5}\left (\sqrt {5}+\sqrt {3}\right )}{5-3} -\dfrac {4\sqrt {3}\left (\sqrt {5}-\sqrt {3}\right )}{5-3}\\ &=10+2\sqrt {15}-\left (2\sqrt {15}-6\right )\\ &=16. \end {align*} Vậy biểu thức đã cho rút gọn thành $16$. \end {enumerate}  \hfill \qedEX 
\end{Solbook}
\begin{Solbook}{3}
 \par \noindent  \begin {enumerate} \item Vì hãng A thu $200$ nghìn đồng tiền mở cửa và $4$ nghìn đồng cho mỗi kilômét nên \[ f(x)=4x+200. \] Vì hãng B thu $80$ nghìn đồng tiền mở cửa và $6$ nghìn đồng cho mỗi kilômét nên \[ g(x)=6x+80. \] \par \item Hai hãng có số tiền thuê bằng nhau khi $f(x)=g(x)$. Ta có \begin {align*} 4x+200&=6x+80\\ 200-80&=6x-4x\\ 120&=2x\\ x&=60. \end {align*} Vậy hai hãng thu cùng số tiền khi quãng đường dài $60$ km. \par Xét hiệu \begin {align*} f(x)-g(x) &=(4x+200)-(6x+80)\\ &=120-2x\\ &=2(60-x). \end {align*} Với $x>60$ thì $60-x<0$, nên $f(x)-g(x)<0$, tức là $f(x)<g(x)$. \par Do đó, nếu quãng đường dài hơn $60$ km thì nên chọn hãng A để tiết kiệm chi phí. \end {enumerate}  \hfill \qedEX 
\end{Solbook}
\begin{Solbook}{4}
 \par \noindent  \begin {enumerate} \item Gọi $x$ (nghìn đồng) là giá niêm yết của nồi cơm điện và $y$ (nghìn đồng) là giá niêm yết của bàn là, với $x>0$, $y>0$. \par Tổng giá niêm yết là $750$ nghìn đồng, nên \[ x+y=750. \tag {1} \] Số tiền được giảm là $20\%$ giá nồi cộng với $10\%$ giá bàn là và bằng $120$ nghìn đồng, nên \[ \dfrac {20}{100}x+\dfrac {10}{100}y=120, \] tức là \[ 0{,}2x+0{,}1y=120. \] Nhân hai vế với $10$, ta được \[ 2x+y=1200. \tag {2} \] Lấy $(2)$ trừ $(1)$: \begin {align*} (2x+y)-(x+y)&=1200-750\\ x&=450. \end {align*} Thay $x=450$ vào $(1)$: \[ 450+y=750 \Rightarrow y=300. \] Cặp $(x;y)=(450;300)$ thỏa điều kiện $x>0$, $y>0$. \par Vậy giá niêm yết của nồi cơm điện là $450\,000$ đồng và giá niêm yết của bàn là là $300\,000$ đồng. \par \item Sau khi giảm giá, số tiền thực tế là \begin {align*} \text {nồi cơm điện: }&450\,000\cdot (100\%-20\%)=450\,000\cdot 0{,}8=360\,000 \text { (đồng)},\\ \text {bàn là: }&300\,000\cdot (100\%-10\%)=300\,000\cdot 0{,}9=270\,000 \text { (đồng)}. \end {align*} Vậy phải trả $360\,000$ đồng cho nồi cơm điện và $270\,000$ đồng cho bàn là. \par \emph {Kiểm tra.} Số tiền giảm là $90\,000+30\,000=120\,000$ (đồng), đúng với giả thiết. \end {enumerate}  \hfill \qedEX 
\end{Solbook}
\begin{Solbook}{5}
 \par \noindent  \begin {enumerate} \item Gọi $M$ là vị trí mắt, $A$ là gốc cột đèn, $B$ là đỉnh cột đèn. Kẻ $MH$ vuông góc với $AB$ tại $H$. Khi đó $MH$ song song với mặt đất, $MH=12$ m và $AH=1{,}5$ m. \par Tam giác $MHB$ vuông tại $H$, $\widehat {HMB}=30^\circ $, cạnh đối của góc này là $HB$ và cạnh kề là $MH$. Do đó \[ \tan 30^\circ =\dfrac {HB}{MH} \Rightarrow HB=MH\cdot \tan 30^\circ =12\cdot \dfrac {\sqrt {3}}{3}=4\sqrt {3}\ \text {(m)}. \] Chiều cao cột đèn là \[ AB=AH+HB=1{,}5+4\sqrt {3}\ \text {(m)}. \] Vì $\sqrt {3}\approx 1{,}732$ nên $4\sqrt {3}\approx 6{,}928$ và \[ AB\approx 1{,}5+6{,}928=8{,}428\approx 8{,}4\ \text {(m)}. \] Vậy chiều cao cột đèn khoảng $8{,}4$ m. \par \item Phần cột đèn nằm phía trên tầm mắt vẫn là $HB=4\sqrt {3}$ m. \par Gọi $d$ (m) là khoảng cách từ học sinh đến gốc cột đèn, $d>0$. Tia nhìn tạo với phương nằm ngang góc $45^\circ $ nên \[ \tan 45^\circ =\dfrac {4\sqrt {3}}{d}. \] Vì $\tan 45^\circ =1$ nên $d=4\sqrt {3}$. \par Vậy học sinh phải đứng cách gốc cột đèn $4\sqrt {3}$ m. \end {enumerate}  \hfill \qedEX 
\end{Solbook}
\begin{Solbook}{6}
 \par \noindent  \immini { \begin {enumerate} \item $A,P,O,M$ cùng nằm trên đường tròn đường kính $OP$ \par Ta có $OP\perp AM, BM\perp AM$ suy ra $OP\parallel BM$. \item $\triangle AOP = \triangle OBN \Rightarrow OP = BN$ \par Ta lại có $BN \parallel OP$ nên $OBNP$ là hình bình hành. \par \item Chứng minh $I$ là trực tâm $\triangle POJ$ suy ra $IJ \perp OP$ (1) \par Chứng minh $K$ là trung điểm $OP$. \par Lại có: $\widehat {APO} = \widehat {OPI} = \widehat {IOP}$ \par Suy ra $\triangle IPO$ cân tại $I$ suy ra $IK \perp OP$ (2) \par Từ (1) và (2) suy ra $I, J, K$ thẳng hàng. \end {enumerate} } { \begin {tikzpicture}[scale=1, font=\footnotesize , line join =round, line cap=round, >=stealth] \def \r {1.8} \path (0,0) coordinate (O) (180:\r ) coordinate (A) (0:\r ) coordinate (B) (-\r ,2.5*\r ) coordinate (x)node[left]{$x$} ($(A)!4/5!(x)$) coordinate (P) ($(O)!(A)!(P)$) coordinate (H) ($(A)!2!(H)$) coordinate (M) ($(O)!1cm!-90:(A)$) coordinate (n') (intersection of B--M and O--n') coordinate (N) (intersection of A--N and O--P) coordinate (K) (intersection of O--N and M--P) coordinate (I) (intersection of O--M and N--P) coordinate (J) ; \draw (O) circle (\r ) (x)--(A)--(B)--(N)--(A)--(M)--(P)--(O)--(N) (P)--(J)--(O) ; \draw [dashed] (K)--(J); \foreach \x /\g in{O/-90,A/180,B/0,P/180,K/160,M/0,N/90,I/120,J/40} \draw [fill=white] (\x ) circle (1.5pt) +(\g :.3) node{$\x $}; \end {tikzpicture} }  \hfill \qedEX 
\end{Solbook}
